% !PRINTED 485
\weberpage{485}

an integral rational function of $t$, whose coefficients are rational
functions of the $\alpha, \alpha_1 \ldots \alpha_{n-1}$, which remain unchanged
when any permutation of the group $P$ is carried out,
because under these permutations the $\omega$ either remain unchanged,
or are permuted among one another in the same way as the $\Theta$. By \S.\,149 b) these coefficients are therefore contained
in $\Omega$. If we then substitute in (5) the value $\Theta$ for the
indeterminate $t$, it follows, since $\varphi'(\Theta)$ is different from zero,
\weq{6}{\omega = \frac{\Phi(\Theta)}{\varphi'(\Theta)}.}

If we apply this to the coefficients of the product
\weq{7}{(u - \alpha)(u - \alpha_1) \ldots (u - \alpha_{r-1}) = \varphi(u, \Theta),}
where $u$ denotes a variable, it follows that this function
of degree $r$, whose roots are the $\alpha, \alpha_1 \ldots \alpha_{r-1}$, can be
expressed rationally in terms of $\Theta$.

Thus the quantity $\alpha$, which is originally the root of an
equation of degree $n$ in $\Omega$, is at the same time the root of an equation
of degree $r$, whose coefficients depend rationally on the root of an
equation of degree $s$. An equation $f(x) = 0$
whose root $\alpha$ has this property is called an \gesperrt{imprimitive equation}.

We can also express what has been proved thus:

\gesperrt{2. The imprimitive field $\Omega(\alpha)$ of degree $n$ passes,
by the adjunction to $\Omega$ of the field of degree $s$, $\Omega' = \Omega(\Theta)$,
into a field $\Omega'(\alpha)$ of degree $r$.}

We wish further to investigate whether the imprimitivity of the
group is a sufficient criterion for the imprimitivity of the
field, that is, whether one can always find imprimitive elements in a field $\Omega(\alpha)$
with imprimitive group.

Let now, therefore, $f(x) = 0$ be an irreducible equation of degree $n$
with imprimitive group, and let (1) be the roots
of this equation, decomposed into the sequences (2) in such a way that the
elements of these sequences are not separated from one another
by the permutations of the group.

We choose any symmetric function $\psi(x, x_1 \ldots x_{r-1})$
in such a way that the values
%CONT>>
